% ============================================================================
% RAPPORT TECHNIQUE COMPLET - CRIST
% Simulateur d'Évaporation Multi-Effet et Cristallisation Industrielle
% FST Settat - Génie des Procédés - 2024/2025
% ============================================================================

\documentclass[12pt,a4paper,twoside]{article}

% Encodage et langue
\usepackage[utf8]{inputenc}
\usepackage[T1]{fontenc}
\usepackage[french]{babel}

% Mathématiques avancées
\usepackage{amsmath,amssymb,amsfonts}
\usepackage{mathtools}

% Mise en page
\usepackage[margin=2.5cm]{geometry}
\usepackage{fancyhdr}
\usepackage{titlesec}

% Graphiques et couleurs
\usepackage{graphicx}
\usepackage{xcolor}

% Tableaux professionnels
\usepackage{booktabs}
\usepackage{array}
\usepackage{multirow}

% Typographie
\usepackage{microtype}
\usepackage{lettrine}
\usepackage{enumitem}

% Références et liens
\usepackage[colorlinks=true,linkcolor=blue!60!black,citecolor=green!50!black,urlcolor=blue!70!black]{hyperref}

% Notes de bas de page
\usepackage[bottom]{footmisc}

% Couleurs personnalisées
\definecolor{cristblue}{RGB}{0,82,147}
\definecolor{cristgreen}{RGB}{39,174,96}

% En-têtes et pieds de page
\pagestyle{fancy}
\fancyhf{}
\fancyhead[LE,RO]{\thepage}
\fancyhead[LO]{\leftmark}
\fancyhead[RE]{CRIST - Rapport Technique}
\renewcommand{\headrulewidth}{0.4pt}

% Style des sections
\titleformat{\section}{\normalfont\Large\bfseries\color{cristblue}}{\thesection}{1em}{}
\titleformat{\subsection}{\normalfont\large\bfseries\color{cristblue!80}}{\thesubsection}{1em}{}
\titleformat{\subsubsection}{\normalfont\normalsize\bfseries}{\thesubsubsection}{1em}{}

% ============================================================================
\begin{document}

% PAGE DE TITRE
\begin{titlepage}
\centering
\vspace*{1.5cm}
{\scshape\LARGE Faculté des Sciences et Techniques\par}
{\scshape\large Settat -- Maroc\par}
\vspace{1.5cm}
\rule{\textwidth}{2pt}\par
\vspace{0.6cm}
{\Huge\bfseries\color{cristblue} CRIST\par}
\vspace{0.4cm}
{\Large\itshape Simulateur d'Évaporation Multi-Effet\\et Cristallisation Industrielle\par}
\vspace{0.4cm}
\rule{\textwidth}{2pt}\par
\vspace{1.5cm}
{\Large\bfseries Rapport Technique Exhaustif\par}
\vspace{0.3cm}
{\large Modélisation Mathématique, Architecture Logicielle,\\Interface Utilisateur et Déploiement Cloud\par}
\vfill
{\large
\begin{tabular}{rl}
\textbf{Filière :} & Génie des Procédés Industriels \\
\textbf{Module :} & Projet de Fin d'Études \\
\textbf{Encadrant :} & Pr. [Nom] \\
\textbf{Année Universitaire :} & 2024 -- 2025 \\
\end{tabular}
}
\vspace{1.5cm}
{\small Document généré le \today}
\end{titlepage}

\tableofcontents
\newpage

% ============================================================================
\section{Introduction Générale}
% ============================================================================

\lettrine[lines=3, findent=3pt, nindent=0pt]{L}{}a concentration des solutions sucrées par évaporation multi-effet et la cristallisation du saccharose constituent des opérations unitaires fondamentales dans l'industrie sucrière mondiale\footnote{Ces procédés représentent environ 40\% de la consommation énergétique totale d'une sucrerie.}. L'optimisation de ces procédés revêt une importance économique majeure, tant en termes de réduction de la consommation énergétique que d'amélioration de la qualité des produits cristallins obtenus.

Le projet CRIST (\textit{CRystallization and Integrated Simulation Tool}) constitue une application web complète permettant la simulation, l'analyse et l'optimisation technico-économique d'un système intégré d'évaporation-cristallisation\footnote{Développé intégralement en Python avec le framework Streamlit et déployé via conteneurisation Docker sur la plateforme Render.}.

\subsection{Contexte Industriel}

L'industrie sucrière marocaine traite annuellement plus de 3 millions de tonnes de betteraves sucrières\footnote{Source : COSUMAR, rapport annuel 2023.}. L'évaporation multi-effet permet de concentrer le jus sucré de 14\% à 65\% de matière sèche, tandis que la cristallisation produit le sucre cristallisé final. L'intégration thermique de ces deux opérations représente un potentiel d'économie d'énergie considérable.

\subsection{Objectifs du Projet}

Les objectifs de ce travail sont multiples :
\begin{enumerate}[label=\roman*., noitemsep]
    \item Développer un modèle mathématique rigoureux de l'évaporation multi-effet basé sur les bilans de matière et d'énergie
    \item Implémenter la résolution de l'équation de bilan de population (PBE\footnote{\textit{Population Balance Equation}.}) pour la cristallisation batch
    \item Réaliser l'analyse technico-économique incluant CAPEX\footnote{\textit{Capital Expenditure} -- Investissement initial.}, OPEX\footnote{\textit{Operating Expenditure} -- Coûts opératoires.} et TAC\footnote{\textit{Total Annualized Cost} -- Coût total annualisé.}
    \item Effectuer l'intégration thermique par la méthode d'analyse pinch
    \item Créer une interface utilisateur web interactive et intuitive
    \item Déployer l'application sur une infrastructure cloud
\end{enumerate}

% ============================================================================
\section{Modélisation de l'Évaporation Multi-Effet}
% ============================================================================

\subsection{Principe Fondamental}

\lettrine[lines=2]{L}{}e système d'évaporation multi-effet exploite le principe de réutilisation de la vapeur générée dans un effet pour chauffer l'effet suivant\footnote{Ce principe thermodynamique permet d'atteindre des économies de vapeur significatives, typiquement 0.8 à 0.9 kg d'eau évaporée par kg de vapeur vive consommée.}. La configuration adoptée est le co-courant\footnote{Configuration où le liquide et la vapeur circulent dans le même sens, de l'effet à haute pression vers l'effet à basse pression. Alternative : contre-courant.}, où le gradient de pression naturel facilite l'écoulement du liquide.

\subsection{Bilans de Matière}

\subsubsection{Bilan Global}

Pour chaque effet $i$ $(i = 1, 2, \ldots, n)$, le bilan de matière global en régime permanent s'écrit :

\begin{equation}
\boxed{F_i = L_i + V_i}
\label{eq:mass_global}
\end{equation}

où les variables sont définies comme suit :
\begin{itemize}[noitemsep]
    \item $F_i$ : débit massique d'alimentation de l'effet $i$ (kg/h)
    \item $L_i$ : débit massique de liqueur concentrée sortante (kg/h)
    \item $V_i$ : débit massique de vapeur d'eau produite (kg/h)
\end{itemize}

\subsubsection{Bilan sur le Soluté}

Le bilan de matière sur le saccharose (soluté non volatil) impose :

\begin{equation}
F_i \cdot x_i = L_i \cdot x_{i+1}
\label{eq:mass_solute}
\end{equation}

où $x$ représente la fraction massique en saccharose. La combinaison des équations \eqref{eq:mass_global} et \eqref{eq:mass_solute} permet d'exprimer le débit de vapeur :

\begin{equation}
\boxed{V_i = F_i \left(1 - \frac{x_i}{x_{i+1}}\right)}
\label{eq:vapor_flow}
\end{equation}

\subsection{Bilans Énergétiques}

\subsubsection{Bilan Enthalpique}

Le bilan énergétique de l'effet $i$ s'établit entre les flux entrant et sortant\footnote{On néglige les pertes thermiques vers l'environnement, hypothèse raisonnable pour des évaporateurs industriels bien calorifugés.} :

\begin{equation}
\boxed{Q_i = F_i \cdot C_{p,i} \cdot (T_{\text{eb},i} - T_{i,\text{in}}) + V_i \cdot \lambda_i}
\label{eq:energy_balance}
\end{equation}

où :
\begin{itemize}[noitemsep]
    \item $Q_i$ : flux thermique transféré dans l'effet $i$ (kW)
    \item $C_{p,i}$ : capacité calorifique massique de la solution (kJ/kg·K)
    \item $T_{\text{eb},i}$ : température d'ébullition dans l'effet $i$ (°C)
    \item $T_{i,\text{in}}$ : température d'entrée du liquide (°C)
    \item $\lambda_i$ : chaleur latente massique de vaporisation (kJ/kg)
\end{itemize}

\subsubsection{Transfert de Chaleur}

La chaleur est fournie par condensation de la vapeur de chauffage :

\begin{equation}
Q_i = V_{i-1} \cdot \lambda_{i-1} \cdot \eta_{\text{th}} = U_i \cdot A_i \cdot \Delta T_i
\label{eq:heat_transfer}
\end{equation}

où $U_i$ est le coefficient global de transfert (W/m²·K), $A_i$ la surface d'échange (m²), et $\Delta T_i$ la force motrice thermique (K).

\subsection{Propriétés Thermophysiques}

\subsubsection{Élévation Ébulliométrique (BPE)}

La présence de soluté dissous élève le point d'ébullition de la solution par rapport à l'eau pure\footnote{Phénomène colligatif décrit par la loi de Raoult modifiée.}. Cette élévation ébulliométrique est calculée par la corrélation empirique\footnote{Coefficients issus de \textit{Beet-Sugar Handbook}, M. Asadi, 2007.} :

\begin{equation}
\boxed{\Delta T_{\text{BPE}} = 0.2253 \, x + 0.0144 \, x^2 + 0.0011 \, x^3}
\label{eq:bpe}
\end{equation}

La température d'ébullition réelle est donc : $T_{\text{eb}} = T_{\text{sat}}(P) + \Delta T_{\text{BPE}}$

\subsubsection{Coefficient de Transfert de Chaleur}

Le coefficient global $U$ dépend de la viscosité de la solution\footnote{La viscosité augmente fortement avec la concentration, réduisant le transfert de chaleur.} :

\begin{equation}
U = U_0 \cdot \left(\frac{\mu_{\text{eau}}}{\mu_{\text{solution}}}\right)^{0.33}
\label{eq:htc}
\end{equation}

avec $U_0 \approx 2500$ W/m²·K pour un évaporateur à circulation forcée.

\subsection{Économie de Vapeur}

L'efficacité énergétique du système multi-effet est caractérisée par l'économie de vapeur :

\begin{equation}
\boxed{\text{SE} = \frac{\sum_{i=1}^{n} V_i}{S} \approx 0.8 \times n}
\label{eq:steam_economy}
\end{equation}

où $S$ est le débit de vapeur vive et $n$ le nombre d'effets. Cette relation empirique\footnote{Valable pour des effets de taille sensiblement égale et des pertes minimales.} montre l'intérêt économique d'augmenter le nombre d'effets.

% ============================================================================
\section{Modélisation de la Cristallisation Batch}
% ============================================================================

\subsection{Fondements Physico-Chimiques}

\lettrine[lines=2]{L}{}a cristallisation par refroidissement est un procédé de séparation basé sur la variation de solubilité avec la température\footnote{Pour le saccharose, la solubilité diminue de 74 g/100g à 80°C à 67 g/100g à 20°C.}. Le contrôle précis de la sursaturation permet d'optimiser la distribution granulométrique des cristaux (CSD\footnote{\textit{Crystal Size Distribution} -- Distribution de taille des cristaux.}).

\subsection{Équilibre Solide-Liquide}

\subsubsection{Courbe de Solubilité}

La solubilité du saccharose dans l'eau est décrite par le polynôme empirique\footnote{Corrélation valide pour $T \in [20, 90]$°C avec une précision de $\pm 0.1$ g/100g.} :

\begin{equation}
\boxed{C^*(T) = 64.397 + 0.07251 \, T + 2.056 \times 10^{-3} \, T^2 - 9.035 \times 10^{-6} \, T^3}
\label{eq:solubility}
\end{equation}

où $C^*$ est exprimée en g de saccharose par 100 g de solution et $T$ en °C.

\subsubsection{Sursaturation}

La sursaturation relative constitue la force motrice de la cristallisation :

\begin{equation}
\boxed{S = \frac{C - C^*(T)}{C^*(T)} = \frac{\Delta C}{C^*}}
\label{eq:supersaturation}
\end{equation}

La zone métastable\footnote{Zone où la solution sursaturée reste stable sans nucléation spontanée.} correspond typiquement à $S < 0.05$ pour le saccharose.

\subsection{Cinétiques de Cristallisation}

\subsubsection{Nucléation Secondaire}

Le taux de nucléation secondaire\footnote{Nucléation induite par les cristaux existants, dominante en cristallisation industrielle.} suit une loi de puissance empirique :

\begin{equation}
\boxed{B = k_b \cdot S^b \cdot m_T^j}
\label{eq:nucleation}
\end{equation}

avec les paramètres typiques pour le saccharose :
\begin{itemize}[noitemsep]
    \item $k_b = 1 \times 10^{12}$ nucléi/m³/s
    \item $b = 2.0$ (ordre par rapport à la sursaturation)
    \item $j = 1.0$ (ordre par rapport à la masse cristalline)
\end{itemize}

\subsubsection{Croissance Cristalline}

La vitesse de croissance linéaire\footnote{Hypothèse de McCabe : croissance indépendante de la taille.} obéit à une loi d'Arrhenius modifiée :

\begin{equation}
\boxed{G = k_g \cdot S^g \cdot \exp\left(-\frac{E_g}{RT}\right)}
\label{eq:growth}
\end{equation}

où $E_g \approx 50$ kJ/mol est l'énergie d'activation de croissance.

\subsection{Équation de Bilan de Population}

L'évolution temporelle de la distribution granulométrique est gouvernée par la PBE\footnote{Équation intégro-différentielle établie par Randolph \& Larson (1971).} :

\begin{equation}
\boxed{\frac{\partial n}{\partial t} + G \frac{\partial n}{\partial L} = B \, \delta(L - L_0) - D(L)}
\label{eq:pbe}
\end{equation}

où $n(L,t)$ est la densité de population (nucléi/m³/m), $L$ la taille caractéristique, et $\delta$ la distribution de Dirac représentant la nucléation à taille $L_0$.

\subsection{Résolution par la Méthode des Moments}

La transformation par les moments\footnote{Cette approche réduit la PBE à un système d'EDO, évitant la discrétisation en taille.} est définie par :

\begin{equation}
\mu_k(t) = \int_0^{\infty} L^k \cdot n(L,t) \, dL
\label{eq:moment_def}
\end{equation}

L'évolution temporelle des quatre premiers moments est :

\begin{equation}
\boxed{\frac{d\mu_k}{dt} = k \cdot G \cdot \mu_{k-1} + B \cdot L_0^k \quad (k = 0, 1, 2, 3)}
\label{eq:moment_ode}
\end{equation}

Les moments donnent accès aux propriétés de la CSD :
\begin{align}
L_{10} &= \frac{\mu_1}{\mu_0} \quad \text{(taille moyenne en nombre)} \\
L_{43} &= \frac{\mu_4}{\mu_3} \quad \text{(diamètre de Sauter)} \\
\text{CV} &= \sqrt{\frac{\mu_2 \mu_0}{\mu_1^2} - 1} \quad \text{(coefficient de variation)}
\end{align}

\subsection{Stratégies de Refroidissement}

\subsubsection{Profil Linéaire}

Le refroidissement linéaire, simple à implémenter, provoque une sursaturation excessive au début du batch\footnote{Conduisant à une nucléation massive et des cristaux fins.} :

\begin{equation}
T(t) = T_0 - \frac{(T_0 - T_f)}{t_{\text{batch}}} \cdot t
\end{equation}

\subsubsection{Profil Optimal de Mullin-Nyvlt}

Le profil cubique\footnote{Conçu pour maintenir la sursaturation approximativement constante.} produit des cristaux plus gros et plus uniformes :

\begin{equation}
\boxed{T(t) = T_f + (T_0 - T_f) \cdot \left(1 - \frac{t}{t_{\text{batch}}}\right)^3}
\label{eq:optimal_cooling}
\end{equation}

% ============================================================================
\section{Dimensionnement des Équipements}
% ============================================================================

\subsection{Surface d'Échange des Évaporateurs}

La surface requise pour chaque effet est calculée par :

\begin{equation}
\boxed{A_i = \frac{Q_i}{U_i \cdot (T_{\text{vapeur},i-1} - T_{\text{eb},i})}}
\label{eq:area_evap}
\end{equation}

\subsection{Volume du Cristallisoir}

Le volume est déterminé par le temps de séjour et la production souhaitée :

\begin{equation}
V_{\text{crist}} = \dot{m}_{\text{cristaux}} \cdot \tau \cdot \frac{1}{\rho_{\text{slurry}} \cdot \phi_s}
\end{equation}

où $\phi_s$ est la fraction volumique de solide ($\approx 0.3$ à 0.4).

\subsection{Surface de Refroidissement}

Le serpentin de refroidissement est dimensionné par :

\begin{equation}
\boxed{A_{\text{serp}} = \frac{Q_{\text{refroid}}}{U_{\text{serp}} \cdot \Delta T_{\text{ml}}}}
\end{equation}

\subsection{Puissance d'Agitation}

La puissance mécanique pour maintenir les cristaux en suspension\footnote{Critère de Zwietering pour la suspension complète.} :

\begin{equation}
P = N_p \cdot \rho_{\text{slurry}} \cdot N^3 \cdot D_a^5
\end{equation}

% ============================================================================
\section{Analyse Économique}
% ============================================================================

\subsection{Coût Total Annualisé (TAC)}

\lettrine[lines=2]{L}{}e TAC représente le critère d'optimisation économique principal\footnote{Permet de comparer des configurations avec différents ratios CAPEX/OPEX.} :

\begin{equation}
\boxed{\text{TAC} = \frac{\text{CAPEX}}{n_{\text{ann}}} + \text{OPEX}}
\label{eq:tac}
\end{equation}

où $n_{\text{ann}}$ est la durée d'amortissement (typiquement 10 ans).

\subsection{Estimation du CAPEX}

\subsubsection{Évaporateurs}

Le coût d'investissement suit une loi d'économie d'échelle\footnote{Exposant de 0.6 typique pour les échangeurs de chaleur.} :

\begin{equation}
C_{\text{évap}} = C_0 \cdot \left(\frac{A}{A_0}\right)^{0.6} \cdot f_{\text{mat}} \cdot f_{\text{pres}} \cdot f_{\text{inst}}
\end{equation}

avec $C_0 = 30\,000$ €  pour $A_0 = 100$ m², $f_{\text{mat}} = 1.3$ (acier inox), $f_{\text{inst}} = 1.5$.

\subsubsection{Cristallisoir}

\begin{equation}
C_{\text{crist}} = 15\,000 \cdot V^{0.7} \quad \text{(€)}
\end{equation}

\subsection{Estimation de l'OPEX}

Les coûts opératoires annuels comprennent :

\begin{align}
\text{OPEX} &= C_{\text{vapeur}} + C_{\text{élec}} + C_{\text{eau}} + C_{\text{maint}} \\
C_{\text{vapeur}} &= \dot{m}_S \cdot c_{\text{vap}} \cdot t_{\text{op}} \quad (c_{\text{vap}} \approx 25 \text{ €/tonne}) \\
C_{\text{maint}} &\approx 0.05 \times \text{CAPEX}
\end{align}

où $t_{\text{op}} \approx 8000$ h/an pour une opération continue.

\subsection{Optimisation du Nombre d'Effets}

L'augmentation de $n$ réduit l'OPEX (moins de vapeur) mais augmente le CAPEX (plus d'équipements). Le nombre optimal minimisant le TAC se situe entre 3 et 6 effets\footnote{Fortement dépendant du ratio prix vapeur/prix équipement.}.

% ============================================================================
\section{Intégration Thermique}
% ============================================================================

\subsection{Analyse Pinch}

La méthode du pincement\footnote{Développée par Linnhoff \textit{et al.} dans les années 1980.} identifie les opportunités de récupération de chaleur entre flux chauds et froids.

\subsubsection{Courbes Composites}

Les courbes composites représentent l'enthalpie cumulée en fonction de la température :

\begin{equation}
H = \int_{T_1}^{T_2} \sum_{\text{flux}} \dot{m}_j \cdot C_{p,j} \, dT
\end{equation}

\subsubsection{Point Pinch et $\Delta T_{\min}$}

Le point pinch correspond au plus faible écart entre courbes :

\begin{equation}
\Delta T_{\min} = \min(T_{\text{chaud}} - T_{\text{froid}})
\end{equation}

Typiquement $\Delta T_{\min} = 10$ K représente un compromis économique\footnote{Valeur plus faible = plus de récupération mais échangeurs plus grands.}.

\subsection{Potentiel de Récupération}

La chaleur récupérable est bornée par :

\begin{equation}
\boxed{Q_{\text{récup}} = \min(Q_{\text{chaud,dispo}}, Q_{\text{froid,requis}})}
\end{equation}

% ============================================================================
\section{Architecture Logicielle}
% ============================================================================

\subsection{Vue d'Ensemble}

\lettrine[lines=2]{L}{}e projet CRIST adopte une architecture en couches\footnote{Respectant le principe de séparation des préoccupations (\textit{Separation of Concerns}).} :

\begin{enumerate}[noitemsep]
    \item \textbf{Couche Présentation} : Interface Streamlit + composants Antd
    \item \textbf{Couche Métier} : Modules de simulation et optimisation
    \item \textbf{Couche Données} : Configuration, export, persistance
\end{enumerate}

\subsection{Structure des Fichiers}

\begin{verbatim}
CRIST/
├── app.py                    # Point d'entrée, authentification
├── config.py                 # Paramètres par défaut
├── pages/                    # Pages Streamlit multipage
│   ├── 1_Home.py            # Vue d'ensemble
│   ├── 2_Evaporator.py      # Simulation MEE
│   ├── 3_Crystallization.py # Simulation cristallisoir
│   ├── 4_Optimization.py    # Optimisation TAC
│   ├── 5_Integration.py     # Analyse pinch
│   └── 6_Results.py         # Export résultats
├── src/                      # Logique métier
│   ├── evaporateurs.py      # Modèle MEE (fsolve)
│   ├── cristallisation.py   # Modèle PBE (odeint)
│   ├── thermodynamique.py   # Propriétés (CoolProp)
│   ├── optimisation.py      # TAC + pinch
│   ├── visualization.py     # Graphiques Plotly
│   ├── export.py            # Excel/PDF
│   └── auth.py              # Authentification
├── .streamlit/secrets.toml   # Credentials (gitignored)
├── Dockerfile               # Image Docker
└── requirements.txt         # Dépendances
\end{verbatim}

\subsection{Modules Principaux}

\begin{description}[style=nextline]
    \item[\texttt{evaporateurs.py}] Implémente les classes \texttt{EvaporatorEffect} et \texttt{MultiEffectEvaporator}. Résolution séquentielle effet par effet via \texttt{scipy.optimize.fsolve}\footnote{Méthode de Newton-Raphson pour systèmes non-linéaires.}.
    
    \item[\texttt{cristallisation.py}] Classe \texttt{BatchCrystallizer} résolvant le système d'EDO des moments par \texttt{scipy.integrate.odeint}\footnote{Intégrateur LSODA adaptatif ordre/pas.}.
    
    \item[\texttt{thermodynamique.py}] Interface avec CoolProp\footnote{Bibliothèque open-source implémentant les équations d'état IAPWS-IF97.} pour les propriétés de l'eau/vapeur : $T_{\text{sat}}(P)$, $\lambda(P)$, $\rho$, $\mu$, $C_p$.
    
    \item[\texttt{visualization.py}] Génération de graphiques Plotly\footnote{Bibliothèque JavaScript avec binding Python pour visualisation interactive.} : profils de température, courbes de sursaturation, CSD, diagrammes pinch.
\end{description}

% ============================================================================
\section{Interface Utilisateur Streamlit}
% ============================================================================

\subsection{Framework Streamlit}

Streamlit\footnote{\url{https://streamlit.io} -- Framework Python pour applications data/ML.} permet de créer des interfaces web interactives sans expertise frontend. Composants utilisés :

\begin{itemize}[noitemsep]
    \item \texttt{st.sidebar} : Navigation latérale
    \item \texttt{st.form} : Formulaires avec validation
    \item \texttt{st.session\_state} : Persistance inter-pages
    \item \texttt{st.plotly\_chart} : Graphiques interactifs
    \item \texttt{st.download\_button} : Export fichiers
    \item \texttt{st.columns}, \texttt{st.expander} : Layout responsive
\end{itemize}

\subsection{Authentification}

L'accès à l'application est protégé par authentification\footnote{Credentials stockés dans \texttt{secrets.toml} (local) ou variables d'environnement (production).} :

\begin{enumerate}[noitemsep]
    \item Formulaire de connexion sur la page principale
    \item Vérification contre \texttt{st.secrets} ou \texttt{os.environ}
    \item État stocké dans \texttt{st.session\_state.authenticated}
    \item Fonction \texttt{require\_auth()} protégeant chaque page
    \item Bouton déconnexion dans le menu latéral
\end{enumerate}

\subsection{Navigation avec streamlit-antd-components}

Le package \texttt{streamlit-antd-components}\footnote{Composants Ant Design portés pour Streamlit.} fournit un menu élégant avec :

\begin{itemize}[noitemsep]
    \item Structure hiérarchique (sous-menus Simulation, Analysis)
    \item Icônes Bootstrap Icons
    \item Navigation via \texttt{st.switch\_page()}
    \item Style professionnel (CSS personnalisé)
\end{itemize}

\subsection{Pages de l'Application}

\begin{table}[htbp]
\centering
\small
\begin{tabular}{lp{7cm}}
\toprule
\textbf{Page} & \textbf{Fonctionnalités} \\
\midrule
Home & Vue d'ensemble, théorie, paramètres par défaut \\
Evaporator & Configuration $n$, $F$, $x$, $P$ ; simulation ; analyse de sensibilité \\
Crystallization & Stratégies de refroidissement ; CSD ; dimensionnement équipement \\
Optimization & Comparaison configurations ; calcul TAC ; recommandation \\
Integration & Courbes composites ; point pinch ; récupération chaleur \\
Results & Dashboard ; export Excel formaté ; génération PDF \\
\bottomrule
\end{tabular}
\end{table}

% ============================================================================
\section{Déploiement Docker et Cloud}
% ============================================================================

\subsection{Conteneurisation Docker}

Docker\footnote{\url{https://docker.com}} encapsule l'application avec ses dépendances pour garantir la reproductibilité\footnote{Élimine le problème ``works on my machine''.}.

\subsubsection{Dockerfile}

\begin{verbatim}
FROM python:3.11-slim

ENV PYTHONDONTWRITEBYTECODE=1
ENV STREAMLIT_SERVER_PORT=10000
ENV STREAMLIT_SERVER_ADDRESS=0.0.0.0

WORKDIR /app
RUN apt-get update && apt-get install -y build-essential
COPY requirements.txt .
RUN pip install --no-cache-dir -r requirements.txt
COPY . .

RUN useradd --create-home appuser && chown -R appuser /app
USER appuser

EXPOSE 10000
CMD ["streamlit", "run", "app.py", "--server.port=10000"]
\end{verbatim}

\subsubsection{Construction et Publication}

\begin{verbatim}
docker build -t anastassaoui/crist-app:latest .
docker login
docker push anastassaoui/crist-app:latest
\end{verbatim}

\subsection{Déploiement sur Render}

Render\footnote{\url{https://render.com} -- Plateforme PaaS supportant Docker.} héberge l'application :

\begin{enumerate}[noitemsep]
    \item Créer un \textit{Web Service} depuis Docker Hub
    \item Image : \texttt{anastassaoui/crist-app:latest}
    \item Variables d'environnement : \texttt{CRIST\_USERNAME}, \texttt{CRIST\_PASSWORD}
    \item Port : 10000
    \item Auto-deploy sur push
\end{enumerate}

\subsection{Gestion des Secrets}

\begin{itemize}[noitemsep]
    \item \textbf{Développement} : \texttt{.streamlit/secrets.toml} (gitignored)
    \item \textbf{Production} : Variables d'environnement Render
    \item \textbf{Fallback} : Credentials par défaut si non configurés
\end{itemize}

% ============================================================================
\section{Dépendances Technologiques}
% ============================================================================

\begin{table}[htbp]
\centering
\small
\begin{tabular}{llp{5.5cm}}
\toprule
\textbf{Package} & \textbf{Version} & \textbf{Rôle} \\
\midrule
streamlit & $\geq 1.28$ & Framework interface web \\
streamlit-antd-components & $\geq 0.3$ & Composants UI Ant Design \\
numpy & $\geq 1.21$ & Calcul matriciel \\
scipy & $\geq 1.7$ & Solveurs (fsolve, odeint) \\
pandas & $\geq 1.3$ & Manipulation DataFrames \\
plotly & $\geq 5.0$ & Visualisation interactive \\
CoolProp & $\geq 6.4$ & Propriétés thermodynamiques \\
thermo & $\geq 0.2$ & Propriétés solutions \\
reportlab & -- & Génération PDF \\
openpyxl & $\geq 3.0$ & Export Excel \\
\bottomrule
\end{tabular}
\end{table}

% ============================================================================
\section{Résultats et Validation}
% ============================================================================

\subsection{Évaporateur 4 Effets}

\begin{table}[htbp]
\centering
\begin{tabular}{lcc}
\toprule
\textbf{Paramètre} & \textbf{Valeur} & \textbf{Unité} \\
\midrule
Débit alimentation & 50\,000 & kg/h \\
Concentration entrée/sortie & 14 / 65 & \% \\
Économie de vapeur & 3.2 & -- \\
Surface totale & 850 & m² \\
Consommation vapeur vive & 12\,500 & kg/h \\
\bottomrule
\end{tabular}
\end{table}

\subsection{Comparaison Stratégies Refroidissement}

\begin{table}[htbp]
\centering
\begin{tabular}{lccc}
\toprule
\textbf{Stratégie} & \textbf{$L_{50}$ (µm)} & \textbf{CV (\%)} & \textbf{Rendement} \\
\midrule
Linéaire & 250 & 45 & 82\% \\
Naturel & 320 & 38 & 85\% \\
Optimal (Mullin-Nyvlt) & 450 & 28 & 91\% \\
\bottomrule
\end{tabular}
\end{table}

% ============================================================================
\section{Conclusion et Perspectives}
% ============================================================================

\lettrine[lines=2]{L}{}e projet CRIST constitue une solution intégrée couvrant l'ensemble de la chaîne :

\begin{itemize}[noitemsep]
    \item Modélisation mathématique rigoureuse (bilans, PBE, moments)
    \item Interface utilisateur moderne (Streamlit + Antd)
    \item Optimisation économique complète (CAPEX, OPEX, TAC)
    \item Intégration thermique (analyse pinch)
    \item Déploiement cloud robuste (Docker + Render)
    \item Authentification sécurisée
\end{itemize}

\textbf{Perspectives :} Extension contre-courant, modèles CFD, optimisation dynamique temps-réel, intégration API PLCs industriels.

% ============================================================================
\section*{Nomenclature}
\addcontentsline{toc}{section}{Nomenclature}

\begin{tabular}{lll}
\toprule
$A$ & Surface d'échange & m² \\
$B$ & Taux de nucléation & nucléi/m³/s \\
$C, C^*$ & Concentration, solubilité & g/100g \\
$C_p$ & Capacité calorifique & kJ/kg·K \\
$G$ & Vitesse de croissance & m/s \\
$L$ & Taille cristalline & m \\
$n$ & Densité de population & nucléi/m³/m \\
$Q$ & Flux thermique & kW \\
$S$ & Sursaturation & -- \\
$T$ & Température & °C \\
$U$ & Coefficient transfert & W/m²·K \\
$V$ & Débit vapeur & kg/h \\
$\lambda$ & Chaleur latente & kJ/kg \\
$\mu_k$ & Moment ordre $k$ & m$^k$/m³ \\
\bottomrule
\end{tabular}

\end{document}
